% chapters/journal/00-abstract.tex
% 摘要

\begin{abstract}
区块链系统中，交易在区块内的执行顺序直接影响区块构建者的手续费收益、用户的确认等待体验以及系统面临的最大可提取价值（MEV）风险。现有排序方法以 Gas 手续费贪心排列为主，在追求短期收益的同时忽略了等待时间均衡性与位置敏感的安全风险；已有强化学习方法多将动作空间定义为在少量预设规则之间切换，无法实现逐笔交易粒度的排序决策。本文将区块构建过程建模为逐笔交易选择的马尔可夫决策过程，采用近端策略优化算法训练排序策略网络，奖励函数采用单步奖励、终止奖励和有效性约束的分层设计，统一建模手续费收益、等待服务、终止公平、位置敏感风险、Gas 利用和非法/提前终止行为，并通过合法性掩码机制保证动作可行性。实验基于统一主场景（$N{=}300$，$p_{\text{risk}}{=}15\%$，$\kappa{=}2.0$）的可配置仿真环境，与七类基线在共享评估池上进行对比，其中含参数基线在固定验证池上调参，正式统计以独立训练 seed 为单位。现阶段结果支持本文方法在收益与综合指标上的折中优势，但公平性与风险指标并非同时最优，约束评估也显示可行率与可行子集收益之间存在权衡。本文的实验结论限定在以太坊生态下的抽象区块构建/L2 Sequencer 排序场景，不直接等价于已完成 L1 PBS/Builder 机制级验证。

\noindent\textbf{关键词：} 区块链；交易排序；深度强化学习；序列决策；区块构建
\end{abstract}
